\documentclass[letter]{article}
\usepackage{color}
\usepackage{graphicx}
\usepackage{pifont}
\usepackage[top=2in, bottom=1.5in, left=1in, right=1in]{geometry}
\usepackage{array,booktabs}
\usepackage{relsize}
\usepackage{xcolor}

\newcolumntype{M}[1]{>{\centering\arraybackslash}m{#1}}

\setlength{\parskip}{0.3cm plus4mm minus3mm}
\setlength{\parindent}{0cm}

\newcommand{\clubs}{{\larger[2]\Pisymbol{psy}{"A7}}}
\newcommand{\spades}{{\larger[2]\Pisymbol{psy}{"AA}}}
\newcommand{\hearts}{{\larger[2]\color{red}\Pisymbol{psy}{"A9}}}
\newcommand{\diamonds}{{\larger[2]\color{red}\Pisymbol{psy}{"A8}}}

\newcommand{\clubsX}[1]{\clubs\,#1}
\newcommand{\spadesX}[1]{\spades\,#1}
\newcommand{\heartsX}[1]{\hearts\,#1}
\newcommand{\diamondsX}[1]{\diamonds\,#1}

\newcommand{\trumpjacks}{\clubsX{J} \spadesX{J} \heartsX{J} \diamondsX{J}}
\newcommand{\nonjackcards}{{\bfseries A 10 K Q 9 8 7}}
\newcommand{\cards}{{\bfseries A K Q J 10 9 8 7}}
\newcommand{\bids}{{\bfseries18, 20, 22, 23, 24, 27, 30, 33, 35, 36, 40, 44,
    45, 46, 48, 50, 54, 55, 59, 60} etc.}

% No abstract title
\renewcommand{\abstractname}{\vspace{-\baselineskip}}

% Most content from, permission to use it as CC-BY-NC 4:
% http://www.pagat.com/schafk/skat.html#gamevalue

\begin{document}

\title{Skat}
  \date{\small{September 27 2015}}
  \maketitle

  \begin{abstract}
    Skat is a three-handed trick taking game. It can also be played by four
    people with the dealer sitting out. This document was created by Tom
    Kazimiers with permission to use content from John McLeod, published on his
    website at \emph{http://www.pagat.com/schafk/skat.html}. It can be
    redistributed under a Creative-Commons Attribution-NonCommericial 4
    license (CC-BY-NC 4).
  \end{abstract}

  \section*{Overview}

  Each player gets 10 of the 32 cards and the remaining two form the
  \emph{skat} and are placed face down in the middle of the table.

  The game starts with an auction of which the winner becomes the
  \emph{declarer}. This player has to play alone against the other two, who play
  as a team. The declarer has the right to exchange the skat cards for two of
  the own cards. The declarer also chooses the trump suit or instead to play
  \emph{Grand} (jacks are the only trump) or \emph{Null} (no trumps and declarer
  must not get any tricks).

  All \emph{card points} combined add up to 120, while some cards don't
  contribute any points. To win, declarer has to to take at least 61 points. For
  the opponents, 60 points is enough to win.

  Each game has a \emph{value}, in \emph{game points}, that depends on the
  trumps or game chosen, the location of the top trumps (\emph{matadors}) and
  whether the declarer used the skat. The declarer generally wins the value of
  the game if successful, and loses the game value (doubled if skat was
  exchanged) if unsuccessful. The game points one wins and loses are separate
  from the card points (which determine whether one wins or loses)


  \section*{Cards}
  
  Skat is often played with German suits. However, French suited cards are
  common as well. There are 32 cards in use. For each suit there is: A K Q J 10
  9 8 7:

  \hfill

  \begin{tabular}{ M{0.5cm} m{2cm} M{0.5cm} m{2cm} m{2cm} }
    {\huge \clubs} & clubs & \includegraphics[height=0.5cm]{images/german_suit_acorn} & acorns \\
    {\huge \spades} & spades & \includegraphics[height=0.5cm]{images/german_suit_gras} & leaves \\
    {\huge \hearts} & hearts & \includegraphics[height=0.5cm]{images/german_suit_herz} & hearts \\
    {\huge \diamonds} & diamonds & \includegraphics[height=0.5cm]{images/german_suit_schell} & bells \\
    {\bfseries \large A} & ace & {\bfseries \large A} & ace \\
    {\bfseries \large K} & king & {\bfseries \large K} & king \\
    {\bfseries \large Q} & queen & {\bfseries \large O} & ober \\
    {\bfseries \large J} & jack & {\bfseries \large U} & unter \\
  \end{tabular}

  \pagebreak

  \subsection*{Ranking of cards and games}
 
  The ranking depends on the game the declarer chooses to play. 

  \subsubsection*{Suit games}

  All four jacks are trump, regardless of the suit. Their ranking follows the
  ranking of suits: \clubs \spades \hearts \diamonds. Next and also trump are the
  remaining seven cards of the chosen suit. This makes eleven trumps in total:
  \trumpjacks{} \nonjackcards. The other three suits contain just seven cards,
  ranked: \nonjackcards.

  \subsubsection*{Grand}

  The four jacks are the only trump: \trumpjacks. The remaining four suits each
  contain seven cards ranking from high to low: \nonjackcards.

  \subsubsection*{Null}

  There are no trumps. The eight cards of each suit rank from high to low:
  \cards. Notice that they are \emph{not} ordered by card value.

  \section*{Deal}

  The first dealer is chosen at random. Thereafter the turn to deal rotates
  clockwise. The dealer shuffles and the player to the dealer's right cuts.
  Cards are dealt face down in four steps: 1. Three cards to each player 2. Two
  cards in the center (the skat) 3. Four cards to each player and finally 4.
  Three cards to each player.

  \subsection*{Card values}

  In suit games and grand, the cards have the following values:

  {\bfseries
    \begin{tabular}{c c c c c c c c c c}
      J &  A & 10 & K & Q & 9 & 8 & 7 \\
      2 & 11 & 10 & 4 & 3 & 0 & 0 & 0 \\
    \end{tabular}
  }

  \section*{Auction}

  Each bid is a number which is the \emph{game value} in \emph{game points} of
  some possible game (see below for calculation of game values). The possible
  bids are therefore \bids. If you bid or accept a bid it means you are prepared
  to play a contract of at least that value in game points.

  The player to the dealer's left is called \emph{forehand (F)}, the player to
  the forehand's left \emph{middlehand (M)}, and the player to the middlehand's
  left is \emph{rearhand (R)}. If there are three players, R is the dealer. If
  there are four players, R is to the dealer's right. Throughout bidding F is
  senior to M who is senior to R. The principle is that a senior player only has
  to \emph{equal} a junior player's bid to win the auction, whereas a junior
  player has to bid \emph{higher} than a senior player to win.

  The first part of the auction takes place between F and M. M speaks first,
  either passing or bidding a number. There is no advantage in making a higher
  than necessary bid so M will normally either pass or begin with the lowest
  bid: 18. If M bids a number, F can either give up the chance to be declarer by
  saying "pass" or compete by saying "yes", which means that F bids the same
  number that M just bid. If F says "yes", M can say "pass", or continue the
  auction with a higher bid, to which F will again answer "yes" or "pass".
  This continues until either F or M drops out of the auction by passing - once
  having passed you get no further opportunity to bid on that hand. 

  The second part of the auction is similar to the first part, but takes place
  between R and the survivor of the first part (i.e. whichever of F and M did
  not pass. As the junior player, R either passes or bids a succession of
  numbers, the first of which must be higher than any number mentioned in the
  first part of the auction. To each number bid by R, the survivor must answer
  "yes" or "pass". The winner of the second part of the auction becomes the
  declarer, and the bid is the last number the declarer said or accepted.

  If both M and R pass without having bid, then F can either be declarer at the
  lowest bid (18), or can throw in the cards without play. If the cards are
  thrown in there is no score for the hand, and the next dealer deals. Here are
  three examples:

  \noindent\fcolorbox{white}{gray!20}{%
    \begin{minipage}[t][5cm]{0.5\textwidth}
      \centering
      \vspace{0pt}
      \begin{tabular}{l l l l}
           F &    M &  R & \\
         \cmidrule(r){1-3}
             &   18 &    & \\
         yes &   20 &    & \\
         yes & pass &    & (F wins first part) \\
             &      & 22 & \\
         yes &      & 23 & \\
         yes &      & 24 & \\
        pass &      &    & (R is declarer in 24) \\
      \end{tabular}
      \vfill
    \end{minipage}
  }
  \begin{minipage}[t][5cm]{0.5\textwidth}
    \centering
    \vspace{0pt}
    \noindent\fcolorbox{white}{gray!20}{%
    \begin{tabular}{l l l l}
        F &    M &  R & \\
       \cmidrule(r){1-3}
          & pass &    & (F wins first part) \\
          &      & 18 & \\
      yes &      & pass & (F is declarer in 18) \\
    \end{tabular}
    }
    \vfill
    \noindent\fcolorbox{white}{gray!20}{%
    \begin{tabular}{l l l l}
         F &    M &  R & \\
       \cmidrule(r){1-3}
           &   18 &    & \\
      pass &      &    & (M wins first part) \\
           &      & 20 & \\
           & pass &    & (R is declarer in 20) \\
    \end{tabular}
    }
  \end{minipage}

  To remember whose turn it is to start the bidding, German players sometimes
  say "geben, hören, sagen" (deal, listen, speak), pointing in turn to dealer,
  forehand and middlehand. If middlehand forgets to begin, forehand can start
  proceedings by saying "I'm forehand" or "I'm listening", or "Speak to me!".

  \section*{Possible Contracts}

  The winner of the bidding is entitled to pick up the two skat cards, add them
  to your hand without showing them to the other players, and discard any two
  cards face down. The cards discarded may include one or both of the cards
  picked up, and their value counts along with your tricks. Having discarded,
  you declare your game. If you looked at the skat, your contract is a skat
  game. There are seven possibilities:

  {\bfseries Diamonds, Hearts, Spades, Clubs} (in which the named suit is trumps and the
   declarer tries to take at least 61 card points),

  {\bfseries Grand} (in which the jacks are the only trumps and the declarer tries to take
   at least 61 card points),

  {\bfseries Null} (in which there are no trumps and the declarer tries to lose every
   trick),

  {\bfseries Null Ouvert} (Open Null) (like Null but with declarer's cards are exposed).

  You may choose not to look at the skat cards, but to play with the 10 cards
  you were originally dealt. If you don't look at the skat you are playing a
  hand game, and again there are seven possibilities: Diamonds Hand, Hearts
  Hand, Spades Hand, Clubs Hand, Grand Hand, Null Hand and Null Ouvert Hand.
  In this case no one must look at the skat cards until after the play.

  If you are declarer in a Suit Hand or Grand Hand game, you can increase the
  value of the game by announcing {\bfseries Schneider} (undertaking to win at
  least 90 card points), or {\bfseries Schwarz} (undertaking to win all the
  tricks), or {\bfseries Open} (Ouvert) (undertaking to win all the tricks with
  your cards exposed). Such announcements must be made before the first trick is
  made. \emph{These announcements are not allowed if declarer has looked at the
  skat}. Also (obviously) they do not apply in Null games.

  \section*{The Play}

  Play is clockwise. No matter who is the declarer, forehand always leads to the
  first trick. Players must follow suit if they can. A player with no card of
  the suit led may play any card. Note that in Suit and Grand games the jacks
  belong to the trump suit, \emph{not} to the suits marked on them. For example
  if hearts are trumps, the jack of clubs is the highest \emph{heart}, and has
  nothing whatever to do with the club suit.

  A trick is won by the highest card of the suit led, unless it contains a
  trump, in which case the highest trump wins it. The winner of a trick leads
  to the next.

  If you are declarer in a {\bfseries Suit} or {\bfseries Grand} game you win if
  the cards in your tricks plus the skat contain at least 61 card points. The
  opponents win if their combined tricks contain at least 60 card points.

  If the declarer's opponents take 30 points or fewer in tricks, they are
  \emph{Schneider}. If they take 31 or more they are said to \emph{be out of
  Scheider}. If they take no tricks at all, they are \emph{Schwarz}. The same
  applies to the declarer - as declarer, you are \emph{Schneider} if you win 30
  card points or less including the skat, and Schwarz if you lose every trick.
  Note that Schwarz depends on tricks not points - if a side wins just one trick
  and it has no card points in it, that is sufficient to get them out of
  Schwarz.

  If you are declarer in {\bfseries Null} or {\bfseries Null Ouvert}, you win the
  game if you manage to lose every trick. If you take a trick, you have lost and
  the play of the hand ceases at that point.

  If you are declarer in an Open (Ouvert) contract - i.e. you are playing
  {\bfseries Null Ouvert} or have announced Open in a Suit or Grand contract -
  you have to spread out your hand face up on the table before the lead to the
  first trick. Play then proceeds normally, and you play from your exposed hand.
  The opponents are not allowed to discuss tactics.

  \section*{Calculating the Value of the Game}

  \subsection*{Suit and Grand contracts}

  The value of a Suit or Grand contract is obtained by multiplying together two
  numbers: the base value and the multiplier. The base value depends on the
  trump suit as follows:

  \begin{tabular}{l c}
    {\bfseries Contract} & {\bfseries Base value} \\
    Diamonds &  9 \\
    Hearts   & 10 \\
    Spades   & 11 \\
    Clubs    & 12 \\
    Grand    & 24 \\
  \end{tabular}

  The multiplier is the sum of all applicable items from the following table
  (n/a = not applicable):

  \begin{tabular}{l c c}
    {\bfseries Multiplier} & {\bfseries Skat game} & {\bfseries Hand game} \\
    Matadors (with or against) & 1 each & 1 each \\
    Game (always applies) & 1 & 1 \\
    Hand (declarer did not look at the skat) & n/a & 1 \\
    Schneider (one side took 90 or more card points) & 1 & 1 \\
    Schneider announced  & n/a & 1 \\
    Schwarz (one side took every trick) & 1 & 1 \\
    Schwarz announced  & n/a & 1 \\
    Open & n/a & 1 \\
  \end{tabular}

  Note that all applicable multipliers count - for example

  \begin{itemize}
    \item if you make the opponents Schwarz, you count the matadors, game,
      Schneider and Schwarz multipliers;
    \item if you announce and make Schneider you count matadors, game, Hand,
      Schneider and Schneider announced;
    \item if you announce and make Schwarz you count matadors, game, Hand,
      Schneider, Schneider announced, Schwarz, Schwarz announced.
  \end{itemize}

  Open contracts are extremely rare: you can only play open if you did not look
  at the skat and you also undertake to win every trick. By implication, an open
  contract includes announcements of Schneider and Schwarz, so you count:
  matadors, game, Hand, Schneider, Schneider announced, Schwarz, Schwarz
  announced, and Open.

  \subsubsection*{Matadors}

  The jack of clubs and any top trumps in unbroken sequence with it are called
  matadors. If as declarer you have such a sequence in your original hand plus
  the skat, you are with that number of matadors. If there is such a sequence in
  the opponents' combined hands, declarer is against that number of matadors.

  {\bfseries Examples of matadors} (Hearts are trumps)

  \begin{tabular}{l l}
    Declarer has & Declarer is \\
    \clubsX{J}, \heartsX{J}, \diamondsX{J}, \heartsX{A}, \heartsX{10},
        \heartsX{Q}, \heartsX{9} & with 1 \\
    \clubsX{J}, \spadesX{J}, \heartsX{J}, \diamondsX{J}, \heartsX{A},
        \heartsX{10}, \heartsX{K} & with 7 \\
    \spadesX{J}, \heartsX{J}, \heartsX{A}, \heartsX{K}, \heartsX{Q},
        \heartsX{7} & against 1 \\
    \diamondsX{J}, \heartsX{A}, \heartsX{10}, \heartsX{K}, \heartsX{Q},
      \heartsX{7} & against 3 \\
  \end{tabular}

  Note that for the purposes of matadors, cards in the skat count as part of
  declarer's hand, even though in a Hand game declarer does not know what is in
  the skat when choosing the game.

  The game multiplier is always counted, whether declarer wins or loses. The
  calculation of the value of a game sounds something like this: "with 2, game
  3, Schneider 4, 4 times spades is 44". The declarer must always be with or
  against at least one matador (the jack of clubs must be somewhere), so the
  smallest possible multiplier is 2, and the smallest possible game value (and
  the lowest possible bid) is 18.

  \subsection*{Null contracts}

  These are easy to score. Each possible Null contract has a fixed value
  unaffected by multipliers. As with all contracts, an unsuccessful declarer
  loses twice the value of the game. The Null values are:

  \begin{tabular}{l c c}
    {\bfseries Contract} & {\bfseries Fixed Value} & {\bfseries Amount lost if
        unsuccessful} \\
    Null & 23 & 46 \\
    Null Hand & 35 & 70 \\
    Null Ouvert & 46 & 92 \\
    Null Ouvert Hand & 59 & 118 \\
  \end{tabular}

  These rather eccentric looking numbers are chosen to fit between the other
  contract values, each being slightly below a multiple of 12. (Before the rule
  change of 1st Jan 1999, Null Hand cost only 35 when lost and Null Ouvert Hand
  cost only 59.)

  \section*{Scoring}

  If declarer wins the game {\bfseries and} the value of the game is as least as
  much as the bid, then the value of the game is added to the declarer's
  cumulative score.

  If the declarer loses the game {\bfseries and} the value of the game is as
  least as much as the bid, then {\bfseries twice} the value of the game is
  subtracted from the declarer's score.

  If the value of the declarer's game turns out to be less than the bid then
  the declarer automatically loses - it does not matter how many card points
  were taken. The amount subtracted from the declarer's score is {\bfseries
  twice} the least multiple of the base value of the game actually played which
  would have fulfilled the bid.

  If as declarer you announce Schneider but take less than 90 card points, or if
  you announce Schwarz or Open and lose a trick, you lose, counting {\bfseries
  all} the multipliers you would have won if you had succeeded.

  {\bfseries Example}: Middlehand holds \heartsX{J}, \diamondsX{J}, \clubsX{10},
  \clubsX{K}, \clubsX{9}, \clubsX{8}, \spades{A}, \diamondsX{A}, \diamondsX{10},
  \diamondsX{7}, and decides to play Clubs Hand. This should normally be worth
  48 game points ("against 2, game 3, hand 4, 4 time clubs is 48"). Rearhand has
  a Null Ouvert and bids up to 46, to which {\bfseries M} says yes. {\bfseries
  M} plays clubs hand and takes 74 card points (including the skat cards), but
  unfortunately the skat contains \clubsX{J}, \hearts{Q}. {\bfseries M} is
  therefore with 1 matador (not against 2 as expected), and the game is worth
  only 36 ("with 1, game 2, hand 3 times clubs"), which is less than the bid.
  {\bfseries M} therefore loses 96 game points (twice the 48 points which would
  be the minimum value in clubs which would fulfill the bid). Had {\bfseries M}
  taken (say) 95 card points, the Schneider multiplier would have increased the
  value of the game to 48 ("with 1, game 2, hand 3, schneider 4 times clubs")
  and {\bfseries M} would have won 48 game points.

  A side effect of the method of scoring is that if there are four players at
  the table, the dealer of a hand is effectively against the declarer, winning
  or losing the same as the declarer's opponents.

  \section*{Variations}

  \subsection*{Kontra and Rekontra}
  
  This variation is very widely played in social games. Either opponent of the
  declarer, at any time before they play their card to the first trick, may say
  {\bfseries kontra}. This doubles the score for the contract, whether won or
  lost. The declarer may immediately answer with {\bfseries rekontra}, which
  doubles the score again.

  Note that it is the {\bfseries score} that is doubled, not the {\bfseries
  value} of the contract. For example suppose I bid up to 20, look at the skat,
  and play in diamonds. I am only with one matador, but am hoping to make the
  opponents schneider. One of the opponents says Kontra, and in the play I win
  85 card points. As I am with 1, the game value is 18, so I have overbid (the
  Kontra does not affect this).  So I lose based on the lowest multiple of
  diamonds which would have been sufficient, namely 27. I lose double because I
  looked at the skat and the score is doubled again for the Kontra, so I lose
  108 game points altogether.

  There is some variation as to when Kontra and Rekontra can be said. Some play
  that Kontra can only be said before the first lead (and a declarer who is
  Forehand must wait before leading to give the opponents an opportunity to
  Kontra).

\end{document}
